\chapter{Validazione del sistema}
\label{cap:validazione}

\notaredazionale{Obiettivo: 30 pagine.}

\section{Benchmark di test}
\notaredazionale{Il dataset diagnostico: 124 casi (\texttt{datasets/eval/qa\_test\_set\_diagnostic.json}), com'è stato
costruito e perché, le categorie diagnostiche (\texttt{single\_fact}, \texttt{multi\_fact}, \texttt{synthesis},
\texttt{conversational\_condensation}, \texttt{image\_grounded}, \texttt{distractor}, \texttt{out\_of\_scope},
\texttt{no\_information}, ecc.) e cosa ciascuna è pensata per stressare. Descrivere anche il disegno deliberato di
documenti ``quasi omonimi'' tra i due corsi indicizzati, usato per testare la disambiguazione.}

\section{Indicatori di accuratezza}
\notaredazionale{Vedi \texttt{REPORT\_ARCHITETTURA.md} \S6.1: le due famiglie di metriche --- RAGAS (\texttt{context\_precision},
\texttt{context\_recall}, \texttt{faithfulness}, \texttt{answer\_relevancy}, \texttt{answer\_correctness},
\texttt{answer\_similarity}) e le metriche deterministiche Top-1/Top-5 --- cosa misura ciascuna, perché servono
entrambe le famiglie (qualità percepita end-to-end vs. qualità isolata del recupero).}

\section{Esperimenti effettuati}
\notaredazionale{Vedi \S6.2--6.3: l'ablation della pipeline di retrieval (13 configurazioni,
\texttt{datasets/eval/retrieval\_screening/}) e il confronto tra modelli generativi
(\texttt{datasets/eval/model\_comparison/}), con relativa discussione dei risultati.

\textbf{Attenzione}: questi risultati sono precedenti ad alcune correzioni successive (fusione dei chunk corti, filtro
di corso diventato bonus di priorità invece di esclusione, decomposizione delle domande multi-parte) --- vanno
rieseguiti per avere numeri aggiornati prima della consegna finale (vedi \S7 di \texttt{REPORT\_ARCHITETTURA.md} per
l'elenco completo dei limiti noti al momento della stesura di questa nota).}

\section{Risultati}
\notaredazionale{Sintesi dei risultati quantitativi (tabelle/grafici), confronto tra configurazioni di retrieval e tra
modelli generativi, discussione del compromesso qualità/latenza.}

\subsection{Esempi di casi errati}
\notaredazionale{Materiale già raccolto durante lo sviluppo e lo hardening del sistema, utile da riportare come casi di
studio concreti (con screenshot/log della risposta sbagliata e della diagnosi):
\begin{enumerate}
    \item una domanda su un importo (marca da bollo) ha ricevuto in risposta un numero corretto ma attribuito al fatto
      sbagliato, per contaminazione tra due blocchi di contesto forniti contemporaneamente al modello;
    \item la pagina ``Iscriversi ad anni successivi al primo'' veniva sistematicamente confusa con la pagina
      sull'Obbligo Formativo Aggiuntivo (OFA), per un chunk introduttivo troppo corto per essere distintivo
      nell'embedding;
    \item una domanda sull'accordo di doppio titolo con la Colombia è stata classificata erroneamente come riferita
      a un corso diverso da quello corretto;
    \item il mancato recupero della pagina ``Cosa dicono i nostri studenti'' su una domanda lunga e lessicalmente
      sovraccarica del nome ufficiale del corso --- limite non ancora risolto al momento della stesura di questa nota,
      utile da riportare anche come lavoro futuro.
\end{enumerate}}
